\subsection{CHANGE OF RING}

PROPOSITION 2. Let A, B be two rings, \(\rho: \mathbf{B} \to \mathbf{A}\) a ring homomorphism and E (resp. F) a right (resp. left) A-module. Then there exists one and only one Z-linear mapping

\[
\phi : \rho_ {*} (\mathbf {E}) \otimes_ {\mathbf {B}} \rho_ {*} (\mathbf {F}) \rightarrow \mathbf {E} \otimes_ {\mathbf {A}} \mathbf {F}\tag{6}
\]

such that, for all \(x \in \mathbf{E}\) and \(y \in \mathbf{F}\) , the image under \(\phi\) of the element \(x \otimes y\) of \(\rho_{*}(\mathbf{E}) \otimes_{\mathbf{B}} \rho_{*}(\mathbf{F})\) is the element \(x \otimes y\) of \(\mathbf{E} \otimes_{\mathbf{A}} \mathbf{F}\) ; this \(\mathbf{Z}\) -linear mapping is surjective.

We consider the mapping \((x,y)\mapsto x\otimes y\) of \(\rho_{\star}(E)\times \rho_{*}(F)\) into \(E\otimes_A F\) ; it is Z-bilinear and, for all \(\beta \in B\) , by definition \((x\rho (\beta))\otimes y = x\otimes (\rho (\beta)y)\) , hence conditions (1) of no. 1 hold, whence the existence and uniqueness of \(\phi\) (no. 1, Proposition 1). The latter assertion follows from the fact that the elements \(x\otimes y\) generate the Z-module \(E\otimes_A F\) .

The mapping (6) is called canonical.

COROLLARY. Let \(\mathfrak{a}\) be a two-sided ideal of A such that \(\mathfrak{a}\) is contained in the annihilator of E and in the annihilator of F, so that E (resp. F) has a canonical left (resp. right) \((A/\mathfrak{a})\)-module structure (§ 1, no. 12). Then the canonical homomorphism (6)

\[
\phi : \mathbf {E} \otimes_ {\mathbf {A}} \mathbf {F} \rightarrow \mathbf {E} \otimes_ {\mathbf {A} / \mathfrak {a}} \mathbf {F}
\]

corresponding to the canonical homomorphism \(\rho: \mathbf{A} \to \mathbf{A} / \mathfrak{a}\) is the identity.

For all \(\bar{\alpha} \in \mathbf{A} / \mathfrak{a}\) , all \(x \in \mathbf{E}\) and all \(y \in \mathbf{F}\) , \(x\bar{\alpha} = x\alpha\) (resp. \(\bar{\alpha}y = \alpha y\) ) for all \(\alpha\) such that \(\rho(\alpha) = \bar{\alpha}\) . If \(\mathbf{C} = \mathbf{Z}^{(\mathbb{E} \times \mathbb{F})}\) , the submodule of \(\mathbf{C}\) generated by the elements \((x\alpha, y) - (x, \alpha y)\) is then equal to the submodule generated by the elements \((x\bar{\alpha}, y) - (x, \bar{\alpha}y)\) .

With the hypotheses and notation of Proposition 2, let \(\mathbf{E}'\) be a right B-module, \(\mathbf{F}'\) a left B-module and consider two semi-linear mappings \(u: \mathbf{E}' \to \mathbf{E}\) , \(v: \mathrm{F}' \to \mathrm{F}\) relative to the homomorphism \(\rho: \mathrm{B} \to \mathrm{A}\) ; \(u\) (resp. \(v\) ) can be considered as a B-linear mapping \(\mathrm{E}' \to \rho_*(\mathrm{E})\) (resp. \(\mathrm{F}' \to \rho_*(\mathrm{F})\) ), whence a Z-linear mapping \(w: \mathrm{E}' \otimes_{\mathrm{B}} \mathrm{F}' \to \rho_*(\mathrm{E}) \otimes_{\mathrm{B}} \rho_*(\mathrm{F})\) such that \(w(x' \otimes y') = u(x') \otimes v(y')\) for \(x' \in \mathrm{E}', y' \in \mathrm{F}'\) ; composing the canonical mapping (6) with this mapping, a Z-linear mapping \(w': \mathrm{E}' \otimes_{\mathrm{B}} \mathrm{F}' \to \mathrm{E} \otimes_{\mathrm{A}} \mathrm{F}\) is hence obtained such that \(w'(x' \otimes y') = u(x') \otimes v(y')\) for \(x' \in \mathrm{E}', y' \in \mathrm{F}'\) ; this is the mapping which will normally be denoted by \(u \otimes v\) if no confusion can arise. Clearly \((u, v) \mapsto u \otimes v\) is a Z-bilinear mapping

\[
\operatorname{Hom} _ {\mathbf {B}} \left(\mathrm{E} ^ {\prime}, \rho_ {*} (\mathrm{E})\right) \times \operatorname{Hom} _ {\mathbf {B}} \left(\mathrm{F} ^ {\prime}, \rho_ {*} (\mathrm{F})\right)\rightarrow \operatorname{Hom} _ {\mathbf {Z}} \left(\mathrm{E} ^ {\prime} \otimes_ {\mathbf {B}} \mathrm{F} ^ {\prime}, \mathrm{E} \otimes_ {\mathbf {A}} \mathrm{F}\right).
\]

Moreover, if C is a third ring, \(\sigma: \mathrm{C} \to \mathrm{B}\) a homomorphism, \(\mathrm{E}''\) a right C-module, \(\mathrm{F}''\) a left C-module, \(u': \mathrm{E}'' \to \mathrm{E}'\) and \(v': \mathrm{F}'' \to \mathrm{F}'\) semi-linear mappings relative to \(\sigma\) , then

\[
(u \circ u ^ {\prime}) \otimes (v \circ v ^ {\prime}) = (u \otimes v) \circ (u ^ {\prime} \otimes v ^ {\prime}).
\]

